\section{Objectives and Endpoints}

This study pairs each objective with a prospectively defined, measurable endpoint
and a scientific justification, presented in the standard three-column form. The
trial is a combined investigation of two regulated articles: perioperative
daraxonrasib (RMC-6236) under the investigational new drug (IND) pathway, and the
on-premises large language model (LLM) directed eight-arm robotic
pancreaticoduodenectomy (Whipple) system under the investigational device
exemption (IDE) pathway. The objective set is therefore organized so that the
drug, the device, and the combined procedure each carry their own primary
safety, dose-finding, and feasibility anchors, with secondary oncologic-surgical
quality measures and exploratory survival and human-factors measures layered
beneath them. Table~1 is the governing objectives-endpoints table; it drives the
analysis plan in \S9 and safety rules in \S6 and \S7.

\noindent \textbf{Abbreviations in Trial Protocol:} DLT, dose-limiting toxicity; ISGPS, International Study
Group on Pancreatic Surgery; LLM, large language model; MTD, maximum tolerated
dose; OS, overall survival; PFS, progression-free survival; RP2D, recommended
Phase 2 dose; SAE, serious adverse event. Human comparator values are drawn from
the 2025 Dutch nationwide cohort of 1000 robotic pancreaticoduodenectomies
\cite{DutchCohort2025} and are descriptive benchmarks only; this single-arm
first-in-human study is not for formal comparison.

\vspace{0.3em}




\noindent\begin{tabularx}{\textwidth}{L{4cm} Y Y}
\toprule
\textbf{Objective} & \textbf{Endpoint} & \textbf{Justification} \\
\midrule
\multicolumn{3}{l}{\textit{Primary}} \\
\midrule
Characterize the early safety of the combined LLM-directed robotic Whipple, perioperative daraxonrasib. & Incidence of device-related or procedure-related serious adverse events (SAEs) through 30 days after the procedure, including Clavien-Dindo grade III or higher complications. & Thirty-day SAE incidence is the standard first-in-human safety anchor and captures the dominant early post-Whipple cascade (fistula, hemorrhage, sepsis, failure to rescue); a device or procedure attribution isolates risk introduced by the investigational system. \\
\addlinespace
Determine the maximum tolerated dose (MTD) and recommended Phase 2 dose (RP2D) of perioperative daraxonrasib. & MTD and RP2D defined by the rate of dose-limiting toxicities (DLTs) within the 28-day DLT window of the 3+3 escalation across dose levels DL1 through DL3. & A 3+3 escalation is the established dose-finding design for cytotoxic and targeted oncology agents and yields a defensible RP2D from a small first-in-human cohort, consistent with the RASolute 302 pan-KRAS program \cite{rasolute302}. \\
\addlinespace
Establish the technical feasibility of the eight-arm robotic Whipple under continuous human oversight. & Proportion of procedures that complete all prespecified assigned operative tasks without unplanned conversion caused by device malfunction or unsafe device behavior. & Task-completion without unsafe conversion is the FDA-recognized early-feasibility readout for a significant-risk device and measures whether the system can perform its intended function within hard cyber-physical limits \cite{fda2013earlyfeasibility}. \\
\midrule
\multicolumn{3}{l}{\textit{Secondary}} \\
\midrule
Estimate the oncologic adequacy of the resection. & R0 (microscopically margin-negative) resection rate determined at day-90 pathology review. & R0 status is the strongest modifiable surgical determinant of survival in resectable PDAC and is the principal oncologic quality measure for the operation. \\
\addlinespace
Estimate the rate of clinically relevant postoperative pancreatic fistula. & International Study Group on Pancreatic Surgery (ISGPS) grade B or C postoperative pancreatic fistula rate \cite{Bassi2017}; human comparator 24.4\% \cite{DutchCohort2025}. & Grade B/C fistula is the leading driver of post-Whipple morbidity and mortality; the 2025 Dutch nationwide robotic cohort provides a contemporaneous human benchmark. \\
\bottomrule  
&&Figure~3 continued, next page.
\end{tabularx}









\noindent\begin{tabularx}{\textwidth}{L{4cm} Y Y}
\toprule
\textbf{Objective} & \textbf{Endpoint} & \textbf{Justification} \\
\midrule
\multicolumn{3}{l}{\textit{Primary}} \\
\midrule
Estimate the rate of major postoperative complications. & Clavien-Dindo grade III or higher complication rate through 90 days; human comparator 41.3\% \cite{DutchCohort2025}. & Clavien-Dindo grade III+ is the validated severity threshold for interventions requiring procedural, endoscopic, or intensive-care management. \\
\addlinespace
Estimate early postoperative mortality. & 30-day and 90-day all-cause mortality; human comparator 3.9\% at 30 days \cite{DutchCohort2025}. & Thirty- and 90-day mortality are the standard short-term survival safety measures for major hepatopancreatobiliary surgery. \\
\addlinespace
Characterize the safety profile of the conversion pathway. & Unplanned conversion rate from the robotic system to open or conventional approach; human comparator 10.1\% \cite{DutchCohort2025}. & Conversion rate contextualizes the feasibility endpoint and benchmarks investigational system against established robot practice. \\
\midrule
\multicolumn{3}{l}{\textit{Exploratory (non-confirmatory)}} \\
\midrule
Describe disease control after the combined intervention. & Progression-free survival (PFS) and overall survival (OS) assessed to 24 months. & PFS and OS provide preliminary, hypothesis-generating signals of oncologic benefit; the small single-arm sample precludes confirmatory inference. \\
\addlinespace
Describe the reliability of the LLM advisory layer. & Concordance between the LLM advisory output and the independent hard-coded safety gate, expressed as a proportion of advisories. & Quantifying advisory-to-gate concordance characterizes the second-opinion oracle and informs autonomy-graduation decisions in later phases. \\
\addlinespace
Describe the fidelity of preoperative simulation to the operative event. & Sim-to-real gap in trajectory (mm) and tip force (newtons) between Phase 0 simulation and the intra-operative record. & The sim-to-real gap validates Phase 0 simulation-validation gate, supports the digital-evidence basis for the device feasibility claim. \\
\bottomrule
\end{tabularx}




\vspace{0.3em}

Figure~3 renders the objective-to-endpoint hierarchy that organizes Table~1,
showing how the three primary anchors (30-day serious adverse event incidence,
daraxonrasib MTD/RP2D, and unsafe-conversion-free task completion) sit above the
secondary oncologic-surgical quality measures and the exploratory survival and
human-factors measures.

\begin{mermaidfig}
\node[mmgoal] (prim) at (0,0)      {PRIMARY\\safety + feasibility};
\node[mmstep] (p1)   at (-4.2,-2.4){Device/procedure\\serious AEs\\through 30 days};
\node[mmstep] (p2)   at (0,-2.4)   {MTD / RP2D of\\daraxonrasib (3+3)};
\node[mmstep] (p3)   at (4.2,-2.4) {Assigned tasks done\\without unsafe\\conversion};
\node[mmgoal] (sec)  at (0,-4.8)   {SECONDARY\\oncologic-surgical};
\node[mmin]   (s1)   at (-4.2,-6.9){R0 rate; ISGPS\\B/C fistula};
\node[mmin]   (s2)   at (4.2,-6.9) {Clavien-Dindo III+;\\30/90-day mortality};
\node[mmgoal] (exp)  at (0,-8.3)   {EXPLORATORY\\non-confirmatory};
\node[mmin]   (e1)   at (0,-10.3)  {PFS/OS to 24 mo;\\LLM concordance;\\sim-to-real gap};
\draw[mmedge] (prim) -- (p1);
\draw[mmedge] (prim) -- (p2);
\draw[mmedge] (prim) -- (p3);
\draw[mmedge] (sec) -- (s1);
\draw[mmedge] (sec) -- (s2);
\draw[mmedge] (exp) -- (e1);
\draw[mmedged] (prim) -- (sec);
\draw[mmedged] (sec) -- (exp);
\end{mermaidfig}
\vspace{-0.7cm}
\figcaption{Figure 3. Objective-to-endpoint hierarchy. Primary safety and
feasibility (30-day serious adverse event incidence, daraxonrasib MTD/RP2D, and
unsafe-conversion-free task completion) sit above secondary oncologic-surgical
quality (R0 rate, ISGPS grade B/C fistula, Clavien-Dindo grade III+, and 30- and
90-day mortality) and exploratory progression-free and overall survival, LLM
advisory concordance, and the sim-to-real gap.}

\subsection{Primary Objectives and Endpoints}

The study carries three co-primary objectives, one for each axis of the combined
IND/IDE investigation: a drug-and-procedure safety objective, a drug dose-finding
objective, and a device technical-feasibility objective.

The \textbf{primary safety objective} is to characterize the early safety of the
combined intervention. Its endpoint is the incidence of device-related or
procedure-related serious adverse events through 30 days after the procedure,
graded for severity under the Clavien-Dindo classification, with grade III or
higher events counted as the principal severe-complication signal. Attribution to
the device or to the procedure is adjudicated by an independent reviewer who is
blinded to the daraxonrasib dose level, so that drug-attributable toxicity and
device-attributable or procedure-attributable harm are separated. The 30-day
window is the standard first-in-human horizon for major hepatopancreatobiliary
surgery and captures the dominant early lethal cascade of pancreatic fistula,
intra-abdominal infection, hemorrhage, sepsis, and failure to rescue.

The \textbf{primary dose-finding objective} is to determine the MTD and RP2D of
perioperative daraxonrasib. Its endpoint is defined operationally by the rate of
dose-limiting toxicities within the 28-day DLT window of a standard 3+3
escalation across dose levels DL1 through DL3 (\S4.3). The MTD is the highest
dose level at which fewer than two of six DLT-evaluable participants experience a
DLT, and the RP2D is selected at or below the MTD with integration of all safety,
tolerability, and available pharmacokinetic data. Daraxonrasib eligibility and
the dose anchor follow the broad pan-KRAS criteria of the RASolute 302 program
\cite{rasolute302}.

The \textbf{primary technical-feasibility objective} is to establish that the
eight-arm robotic Whipple can perform its intended function under continuous human
oversight. Its endpoint is the proportion of procedures that complete all
prespecified assigned operative tasks (dissection of the superior mesenteric and
portal veins, hepatic-artery control, the three anastomoses, hemostasis, and
drain placement) without an unplanned conversion caused by device malfunction or
unsafe device behavior. A protocol-defined conversion undertaken for patient
safety at discretion of supervising surgeon, in the absence of device
malfunction, does not count against this endpoint but is recorded and reported.
Task-completion without unsafe conversion is early-feasibility readout
recognized by FDA for a significant-risk device \cite{fda2013earlyfeasibility}.

\subsection{Secondary Objectives and Endpoints}

The secondary objectives characterize the oncologic adequacy and the
surgical-quality profile of the combined intervention against the contemporaneous
human benchmark. Each secondary endpoint is summarized descriptively with an
exact binomial confidence interval; the study is not powered for formal hypothesis
testing against the comparators.

The \textbf{R0 resection rate} is the proportion of participants with a
microscopically margin-negative resection at day-90 pathology review, the single
strongest modifiable surgical determinant of survival in resectable PDAC. The
\textbf{ISGPS grade B/C postoperative pancreatic fistula rate} \cite{Bassi2017}
is the leading driver of post-Whipple morbidity; the 2025 Dutch nationwide
robotic cohort reported a grade B/C fistula rate of 24.4\% \cite{DutchCohort2025}.
The \textbf{Clavien-Dindo grade III or higher complication rate} through 90 days
captures major complications requiring procedural, endoscopic, or intensive-care
management; the Dutch comparator is 41.3\% \cite{DutchCohort2025}.
\textbf{Thirty-day and 90-day all-cause mortality} are the short-term survival
safety measures, with a Dutch 30-day comparator of 3.9\% \cite{DutchCohort2025}.
The \textbf{unplanned conversion rate} from the robotic system to an open or
conventional approach contextualizes the feasibility endpoint; the Dutch
comparator is 10.1\% \cite{DutchCohort2025}. Taken together, these comparators
provide a transparent reference frame, but the small single-arm sample means they
are interpreted as benchmarks rather than as targets for inferential comparison.

\subsection{Exploratory Objectives and Endpoints}

The exploratory objectives generate hypotheses for later-phase study and are
explicitly non-confirmatory; no exploratory analysis is used to make a
go or no-go decision on the intervention and all are reported with that caveat.

\textbf{Progression-free survival and overall survival} are assessed to 24 months
from the day of surgery using Response Evaluation Criteria in Solid Tumors
imaging at day 90 and every 12 weeks thereafter, providing a preliminary signal of
oncologic benefit that the sample size cannot confirm. \textbf{LLM advisory
concordance} is the proportion of LLM advisory outputs that agree with the
independent hard-coded safety gate, quantifying the reliability of the
second-opinion oracle and informing autonomy-graduation decisions reserved for
later phases (\S4.1). The \textbf{sim-to-real gap} is the difference between the
Phase 0 simulation prediction and the intra-operative record, reported as the
trajectory gap in millimeters (target less than 2 mm) and the tip-force gap in
newtons (target less than 0.5 N); it validates the fidelity of the
simulation-validation gate and underpins the digital-evidence basis for the
device feasibility claim. Because these measures are descriptive, they are
summarized with point estimates and dispersion only, without confirmatory
significance testing.
